\documentclass[12pt]{article}

% Idioma y codificación
\usepackage[spanish, es-tabla, es-notilde]{babel}       %es-tabla para que se titule "Tabla"
\usepackage[utf8]{inputenc}

% Márgenes
\usepackage[a4paper,top=3cm,bottom=2.5cm,left=3cm,right=3cm]{geometry}

% Comentarios de bloque
\usepackage{verbatim}

% Paquetes de links
\usepackage[hidelinks]{hyperref}    % Permite enlaces
\usepackage{url}                    % redirecciona a la web

% Más opciones para enumeraciones
\usepackage{enumitem}

% Personalizar la portada
\usepackage{titling}

% Paquetes de tablas
\usepackage{multirow}

% Para añadir el símbolo de euro
\usepackage{eurosym}

% Paquete para subrayar con diferentes estilos
\usepackage[normalem]{ulem}


%------------------------------------------------------------------------

%Paquetes de figuras
\usepackage{caption}
\usepackage{subcaption} % Figuras al lado de otras
\usepackage{float}      % Poner figuras en el sitio indicado H.


% Paquetes de imágenes
\usepackage{graphicx}       % Paquete para añadir imágenes
\usepackage{transparent}    % Para manejar la opacidad de las figuras

% Paquete para usar colores
\usepackage[dvipsnames, table, xcdraw]{xcolor}
\usepackage{pagecolor}      % Para cambiar el color de la página

% Habilita tamaños de fuente mayores
\usepackage{fix-cm}

% Para los gráficos
\usepackage{tikz}
\usepackage{forest}

% Para poder situar los nodos en los grafos
\usetikzlibrary{positioning}

% Para evitar problemas entre babel y tikz
\usetikzlibrary{babel}


%------------------------------------------------------------------------

% Paquetes de matemáticas
\usepackage{mathtools, amsfonts, amssymb, mathrsfs}
\usepackage[makeroom]{cancel}     % Simplificar tachando
\usepackage{polynom}    % Divisiones y Ruffini
\usepackage{units} % Para poner fracciones diagonales con \nicefrac

\usepackage{pgfplots}   %Representar funciones
\pgfplotsset{compat=1.18}  % Versión 1.18

\usepackage{tikz-cd}    % Para usar diagramas de composiciones
\usetikzlibrary{calc}   % Para usar cálculo de coordenadas en tikz

%Definición de teoremas, etc.
\usepackage{amsthm}
%\swapnumbers   % Intercambia la posición del texto y de la numeración

\theoremstyle{plain}

\makeatletter
\@ifclassloaded{article}{
  % --- Configuración para ARTICLE ---
  \newtheorem{teo}{Teorema}[section]
  \theoremstyle{definition}
  \newtheorem{definicion}{Definición}[section]
  \newtheorem{ejercicio}{Ejercicio}[section]

  \renewcommand{\theteo}{%
    \ifnum\value{section}=0 \arabic{teo}\else \thesection.\arabic{teo}\fi
  }
  \renewcommand{\thedefinicion}{%
    \ifnum\value{section}=0 \arabic{definicion}\else \thesection.\arabic{definicion}\fi
  }
  \renewcommand{\theejercicio}{%
    \ifnum\value{section}=0 \arabic{ejercicio}\else \thesection.\arabic{ejercicio}\fi
  }
}{
  % --- Configuración para BOOK / REPORT ---
  \newtheorem{teo}{Teorema}[chapter]
  \theoremstyle{definition}
  \newtheorem{definicion}{Definición}[chapter]
  \newtheorem{ejercicio}{Ejercicio}[section]

  % Si no hay sección: Capítulo.Teorema (ej. 2.1)
  % Si hay sección:    Capítulo.Sección.Teorema (ej. 2.3.1)
  \renewcommand{\theteo}{%
    \ifnum\value{section}=0
      \thechapter.\arabic{teo}%
    \else
      \thesection.\arabic{teo}% (\thesection ya contiene el número de capítulo)
    \fi
  }
  \renewcommand{\thedefinicion}{%
    \ifnum\value{section}=0
      \thechapter.\arabic{definicion}%
    \else
      \thesection.\arabic{definicion}%
    \fi
  }
  \renewcommand{\theejercicio}{%
    \ifnum\value{section}=0
      \arabic{ejercicio}%
    \else
      \thesection.\arabic{ejercicio}% Formato por defecto: ej. 2.2.1
    \fi
  }
}
\makeatother

% Teoremas dependientes o con contador compartido
\newtheorem{coro}{Corolario}[teo]
\newtheorem{prop}[teo]{Proposición}
\newtheorem{lema}[teo]{Lema}
\newtheorem*{lema*}{Lema}

\theoremstyle{remark}
\newtheorem*{observacion}{Observación}

\theoremstyle{definition}
\newtheorem*{notacion}{Notación}
\newtheorem*{ejemplo}{Ejemplo}
\newtheorem*{ejercicio*}{Ejercicio}


% \renewcommand\qedsymbol{$\blacksquare$}         % Cambiar símbolo QED
%------------------------------------------------------------------------

% Paquetes para encabezados
\usepackage{fancyhdr}
\pagestyle{fancy}
\fancyhf{}

\newcommand{\helv}{ % Modificación tamaño de letra
\fontfamily{}\fontsize{12}{12}\selectfont}
\setlength{\headheight}{15pt} % Amplía el tamaño del índice


%\usepackage{lastpage}   % Referenciar última pag   \pageref{LastPage}
%\fancyfoot[C]{%
%  \begin{minipage}{\textwidth}
%    \centering
%    ~\\
%    \thepage\\
%    \href{https://losdeldgiim.github.io/}{\texttt{\footnotesize losdeldgiim.github.io}}
%  \end{minipage}
%}
\fancyfoot[C]{\thepage}
\fancyfoot[R]{\href{https://losdeldgiim.github.io/}{\texttt{\footnotesize losdeldgiim.github.io}}}

%------------------------------------------------------------------------

% Conseguir que no ponga "Capítulo 1". Sino solo "1."
\makeatletter
\@ifclassloaded{book}{
  \renewcommand{\chaptermark}[1]{\markboth{\thechapter.\ #1}{}} % En el encabezado
    
  \renewcommand{\@makechapterhead}[1]{%
  \vspace*{50\p@}%
  {\parindent \z@ \raggedright \normalfont
    \ifnum \c@secnumdepth >\m@ne
      \huge\bfseries \thechapter.\hspace{1em}\ignorespaces
    \fi
    \interlinepenalty\@M
    \Huge \bfseries #1\par\nobreak
    \vskip 40\p@
  }}
}
\makeatother

%------------------------------------------------------------------------
% Paquetes de cógido
\usepackage{minted}
\renewcommand\listingscaption{Código fuente}

\usepackage{fancyvrb}
% Personaliza el tamaño de los números de línea
\renewcommand{\theFancyVerbLine}{\small\arabic{FancyVerbLine}}

% Estilo para C++
\newminted{cpp}{
    frame=lines,
    framesep=2mm,
    baselinestretch=1.2,
    linenos,
    escapeinside=||
}

% para minted
\definecolor{LightGray}{rgb}{0.95,0.95,0.92}
\setminted{
    linenos=true,
    stepnumber=5,
    numberfirstline=true,
    autogobble,
    breaklines=true,
    breakautoindent=true,
    breaksymbolleft=,
    breaksymbolright=,
    breaksymbolindentleft=0pt,
    breaksymbolindentright=0pt,
    breaksymbolsepleft=0pt,
    breaksymbolsepright=0pt,
    fontsize=\footnotesize,
    bgcolor=LightGray,
    numbersep=10pt
}


\usepackage{listings} % Para incluir código desde un archivo

\renewcommand\lstlistingname{Código Fuente}
\renewcommand\lstlistlistingname{Índice de Códigos Fuente}

% Definir colores
\definecolor{vscodepurple}{rgb}{0.5,0,0.5}
\definecolor{vscodeblue}{rgb}{0,0,0.8}
\definecolor{vscodegreen}{rgb}{0,0.5,0}
\definecolor{vscodegray}{rgb}{0.5,0.5,0.5}
\definecolor{vscodebackground}{rgb}{0.97,0.97,0.97}
\definecolor{vscodelightgray}{rgb}{0.9,0.9,0.9}

% Configuración para el estilo de C similar a VSCode
\lstdefinestyle{vscode_C}{
  backgroundcolor=\color{vscodebackground},
  commentstyle=\color{vscodegreen},
  keywordstyle=\color{vscodeblue},
  numberstyle=\tiny\color{vscodegray},
  stringstyle=\color{vscodepurple},
  basicstyle=\scriptsize\ttfamily,
  breakatwhitespace=false,
  breaklines=true,
  captionpos=b,
  keepspaces=true,
  numbers=left,
  numbersep=5pt,
  showspaces=false,
  showstringspaces=false,
  showtabs=false,
  tabsize=2,
  frame=tb,
  framerule=0pt,
  aboveskip=10pt,
  belowskip=10pt,
  xleftmargin=10pt,
  xrightmargin=10pt,
  framexleftmargin=10pt,
  framexrightmargin=10pt,
  framesep=0pt,
  rulecolor=\color{vscodelightgray},
  backgroundcolor=\color{vscodebackground},
}

%------------------------------------------------------------------------
% La etiqueda IA sirve para indicar a partir de donde le pasamos el preambulo a la IA para el escaneo
% IA

% Comandos definidos
\newcommand{\bb}[1]{\mathbb{#1}}
\newcommand{\cc}[1]{\mathcal{#1}}

% I prefer the slanted \leq
\let\oldleq\leq % save them in case they're every wanted
\let\oldgeq\geq
\renewcommand{\leq}{\leqslant}
\renewcommand{\geq}{\geqslant}

% Si y solo si
\newcommand{\sii}{\iff}

% MCD y MCM
\DeclareMathOperator{\mcd}{mcd}
\DeclareMathOperator{\mcm}{mcm}

% Signo
\DeclareMathOperator{\sgn}{sgn}

% Imagen
\DeclareMathOperator{\Ima}{Im}

% Aplicación identidad
\newcommand{\id}{\mathrm{Id}}

% Letras griegas
\newcommand{\eps}{\epsilon}
\newcommand{\veps}{\varepsilon}
\newcommand{\lm}{\lambda}

\newcommand{\ol}{\overline}
\newcommand{\ul}{\underline}
\newcommand{\wt}{\widetilde}
\newcommand{\wh}{\widehat}

\let\oldvec\vec
\renewcommand{\vec}{\overrightarrow}

% Derivadas parciales
\newcommand{\del}[2]{\frac{\partial #1}{\partial #2}}
\newcommand{\Del}[3]{\frac{\partial^{#1} #2}{\partial #3^{#1}}}
\newcommand{\deld}[2]{\dfrac{\partial #1}{\partial #2}}
\newcommand{\Deld}[3]{\dfrac{\partial^{#1} #2}{\partial #3^{#1}}}


\newcommand{\AstIg}{\stackrel{(\ast)}{=}}
\newcommand{\Hop}{\stackrel{L'H\hat{o}pital}{=}}

\newcommand{\red}[1]{{\color{red}#1}} % Para integrales, destacar los cambios.

% Método de integración
\newcommand{\MetInt}[2]{
    \left[\begin{array}{c}
        #1 \\ #2
    \end{array}\right]
}

% Declarar aplicaciones
% 1. Nombre aplicación
% 2. Dominio
% 3. Codominio
% 4. Variable
% 5. Imagen de la variable
\newcommand{\Func}[5]{
    \begin{equation*}
        \begin{array}{rrll}
            \displaystyle #1:& \displaystyle  #2 & \longrightarrow & \displaystyle  #3\\
               & \displaystyle  #4 & \longmapsto & \displaystyle  #5
        \end{array}
    \end{equation*}
}

%------------------------------------------------------------------------


\newcommand{\myhref}[2]{\href{#1}{\color{blue}\uline{#2}}}


\begin{document}

    % 1. Foto de fondo
    % 2. Título
    % 3. Encabezado Izquierdo
    % 4. Color de fondo
    % 5. Coord x del titulo
    % 6. Coord y del titulo
    % 7. Fecha

    
    % 1. Foto de fondo
% 2. Título
% 3. Encabezado Izquierdo
% 4. Color de fondo
% 5. Coord x del titulo
% 6. Coord y del titulo
% 7. Fecha
% 8. Autor

\newcommand{\portada}[8]{

    \portadaBase{#1}{#2}{#3}{#4}{#5}{#6}{#7}{#8}
    \portadaBook{#1}{#2}{#3}{#4}{#5}{#6}{#7}{#8}
}

\newcommand{\portadaExamen}[8]{

    \portadaBase{#1}{#2}{#3}{#4}{#5}{#6}{#7}{#8}
    \portadaArticle{#1}{#2}{#3}{#4}{#5}{#6}{#7}{#8}
}




\newcommand{\portadaBase}[8]{

    % Tiene la portada principal y la licencia Creative Commons
    
    % 1. Foto de fondo
    % 2. Título
    % 3. Encabezado Izquierdo
    % 4. Color de fondo
    % 5. Coord x del titulo
    % 6. Coord y del titulo
    % 7. Fecha
    % 8. Autor
    
    
    \thispagestyle{empty}               % Sin encabezado ni pie de página
    \newgeometry{margin=0cm}        % Márgenes nulos para la primera página
    
    
    % Encabezado
    \fancyhead[L]{\helv #3}
    \fancyhead[R]{\helv \nouppercase{\leftmark}}
    
    
    \pagecolor{#4}        % Color de fondo para la portada
    
    \begin{figure}[p]
        \centering
        \transparent{0.3}           % Opacidad del 30% para la imagen
        
        \includegraphics[width=\paperwidth, keepaspectratio]{../../_assets/#1}
    
        \begin{tikzpicture}[remember picture, overlay]
            \node[anchor=north west, text=white, opacity=1, font=\fontsize{60}{90}\selectfont\bfseries\sffamily, align=left] at (#5, #6) {#2};
            
            \node[anchor=south east, text=white, opacity=1, font=\fontsize{12}{18}\selectfont\sffamily, align=right] at (9.7, 3) {\href{https://losdeldgiim.github.io/}{\textbf{Los Del DGIIM}, \texttt{\footnotesize losdeldgiim.github.io}}};
            
            \node[anchor=south east, text=white, opacity=1, font=\fontsize{12}{15}\selectfont\sffamily, align=right] at (9.7, 1.8) {Doble Grado en Ingeniería Informática y Matemáticas\\Universidad de Granada};
        \end{tikzpicture}
    \end{figure}
    
    
    \restoregeometry        % Restaurar márgenes normales para las páginas subsiguientes
    \pagecolor{white}       % Restaurar el color de página
    
    
    \newpage
    \thispagestyle{empty}               % Sin encabezado ni pie de página
    \begin{tikzpicture}[remember picture, overlay]
        \node[anchor=south west, inner sep=3cm] at (current page.south west) {
            \begin{minipage}{0.5\paperwidth}
                \href{https://creativecommons.org/licenses/by-nc-nd/4.0/}{
                    \includegraphics[height=2cm]{../../_assets/Licencia.png}
                }\vspace{1cm}\\
                Esta obra está bajo una
                \href{https://creativecommons.org/licenses/by-nc-nd/4.0/}{
                    Licencia Creative Commons Atribución-NoComercial-SinDerivadas 4.0 Internacional (CC BY-NC-ND 4.0).
                }\\
    
                Eres libre de compartir y redistribuir el contenido de esta obra en cualquier medio o formato, siempre y cuando des el crédito adecuado a los autores originales y no persigas fines comerciales. 
            \end{minipage}
        };
    \end{tikzpicture}
    
    
    
    % 1. Foto de fondo
    % 2. Título
    % 3. Encabezado Izquierdo
    % 4. Color de fondo
    % 5. Coord x del titulo
    % 6. Coord y del titulo
    % 7. Fecha
    % 8. Autor


}


\newcommand{\portadaBook}[8]{

    % 1. Foto de fondo
    % 2. Título
    % 3. Encabezado Izquierdo
    % 4. Color de fondo
    % 5. Coord x del titulo
    % 6. Coord y del titulo
    % 7. Fecha
    % 8. Autor

    % Personaliza el formato del título
    \pretitle{\begin{center}\bfseries\fontsize{42}{56}\selectfont}
    \posttitle{\par\end{center}\vspace{2em}}
    
    % Personaliza el formato del autor
    \preauthor{\begin{center}\Large}
    \postauthor{\par\end{center}\vfill}
    
    % Personaliza el formato de la fecha
    \predate{\begin{center}\huge}
    \postdate{\par\end{center}\vspace{2em}}
    
    \title{#2}
    \author{\href{https://losdeldgiim.github.io/}{Los Del DGIIM, \texttt{\large losdeldgiim.github.io}}
    \\ \vspace{0.5cm}#8}
    \date{Granada, #7}
    \maketitle
    
    \tableofcontents
}




\newcommand{\portadaArticle}[8]{

    % 1. Foto de fondo
    % 2. Título
    % 3. Encabezado Izquierdo
    % 4. Color de fondo
    % 5. Coord x del titulo
    % 6. Coord y del titulo
    % 7. Fecha
    % 8. Autor

    % Personaliza el formato del título
    \pretitle{\begin{center}\bfseries\fontsize{42}{56}\selectfont}
    \posttitle{\par\end{center}\vspace{2em}}
    
    % Personaliza el formato del autor
    \preauthor{\begin{center}\Large}
    \postauthor{\par\end{center}\vspace{3em}}
    
    % Personaliza el formato de la fecha
    \predate{\begin{center}\huge}
    \postdate{\par\end{center}\vspace{5em}}
    
    \title{#2}
    \author{\href{https://losdeldgiim.github.io/}{Los Del DGIIM, \texttt{\large losdeldgiim.github.io}}
    \\ \vspace{0.5cm}#8}
    \date{Granada, #7}
    \thispagestyle{empty}               % Sin encabezado ni pie de página
    \maketitle
    \vfill
}
    \portadaExamen{ffccA4.jpg}{Topología II\\Examen I}{Topología II. Examen I}{MidnightBlue}{-8}{28}{2025}{José Juan Urrutia Milán}

    \begin{description}
        \item[Asignatura] Topología II.
        \item[Curso Académico] 2023/24.
        \item[Grado] Grado en Matemáticas.
        \item[Grupo] Grupo A.
        % \item[Profesor] ---.
        \item[Descripción] Prueba del Tema 1.
        \item[Fecha] 25 de abril de 2024.
        % \item[Duración] ---.
    \end{description}
    \newpage


    % ------------------------------------
    
    \begin{ejercicio}[7.5 puntos]
        Sea $X_n = \{x\in \mathbb{R}^{n+1} : \|x\| \geq 2\}$ con la topología usual inducida.
        \begin{enumerate}[label=\alph*)]
            \item Encontrar un retracto de deformación de $X_n$ distinto de $X_n$ (detallar la retracción y la homotopía).
            \item Estudiar si $X_1 \cong X_n$ cuando $n\geq 2$.
        \end{enumerate}
    \end{ejercicio}

    \begin{ejercicio}[2.5 puntos]
        Sea $X$ un conjunto con $|X|\geq 2$. Fijado $x_0\in X$ se considera la topología
        \begin{equation*}
            T = \{U\subseteq X : x_0\notin U\}\cup \{X\}
        \end{equation*}
        ¿Es $(X,T)$ contráctil?
    \end{ejercicio}

    \newpage
    \setcounter{ejercicio}{0} % Reiniciar contador de ejercicios
    \noindent
    \textbf{Solución.}

    \begin{ejercicio}[7.5 puntos]
        Sea $X_n = \{x\in \mathbb{R}^{n+1} : \|x\| \geq 2\}$ con la topología usual inducida.
        \begin{enumerate}[label=\alph*)]
            \item Encontrar un retracto de deformación de $X_n$ distinto de $X_n$ (detallar la retracción y la homotopía).

                Consideramos el conjunto:
                \begin{equation*}
                    S_n = \{x\in \mathbb{R}^{n+1} : \|x\| = 2\}
                \end{equation*}

                Antes de demostrar formalmente que, efectivamente, $S_n$ es un retracto de deformación de $X_n$, veamos cómo debe ser el retracto. La idea es contraer cada punto $x\in X_n$ hacia su punto correspondiente en la esfera $S_n$ a lo largo de la recta que une el origen con $x$, como se puede observar en la Figura~\ref{fig:retracto_deformacion}.
                \begin{figure}[ht]
                    \centering
                    \begin{tikzpicture}
                        % Región X_1: exterior del círculo de radio 2
                        \fill[blue!8] (-3.2,-3.2) rectangle (3.2,3.2);
                        % Agujero abierto ||x|| < 2
                        \fill[white] (0,0) circle (1.6cm);

                        % Ejes coordenados
                        \draw[->, thin, gray!70] (-3.4,0) -- (3.4,0) node[right, black] {$x_1$};
                        \draw[->, thin, gray!70] (0,-3.4) -- (0,3.4) node[above, black] {$x_2$};

                        % Circunferencia S_1 (radio 2 representado a escala como 1.6cm)
                        \draw[thick, blue!80!black] (0,0) circle (1.6cm);
                        \node[blue!80!black, font=\small] at (135:1.95cm) {$S_1$};
                        \node[blue!60!black, font=\small] at (-2.5, 2.5) {$X_1$};

                        % Radio 2 indicado
                        \draw[dashed, gray] (0,0) -- (-15:1.6cm) node[midway, below, font=\scriptsize, black] {$2$};

                        % Bucle con varios puntos (ángulo / radio)
                        \foreach \ang/\rad in {25/3.0, 72/2.3, 118/2.8, 165/2.1, 210/3.1, 255/2.4, 305/2.9} {
                            \draw[dotted, gray!60] (0,0) -- (\ang:1.6cm);
                            \draw[->, thick, red!75!black] (\ang:\rad cm) -- (\ang:1.6cm);
                            \fill[red!75!black] (\ang:\rad cm) circle (1.7pt);
                            \fill[blue!80!black] (\ang:1.6cm) circle (1.7pt);
                        }

                        % Punto destacado con etiquetas explícitas
                        \def\angMain{42}
                        \def\radMain{2.85}
                        \draw[dotted, thick, gray!80] (0,0) -- (\angMain:1.6cm);
                        \draw[->, very thick, red!75!black] (\angMain:\radMain cm) -- (\angMain:1.6cm);
                        \fill[red!75!black] (\angMain:\radMain cm) circle (2.2pt) node[above right, font=\small] {$x$};

                        % Origen
                        \fill[black] (0,0) circle (1.5pt) node[below left, font=\scriptsize] {$0$};
                    \end{tikzpicture}
                    \caption{Representación de la retracción por deformación radial de $X_1$ sobre $S_1$ en $\mathbb{R}^2$.}
                    \label{fig:retracto_deformacion}
                \end{figure}

                Para definir la homotopía que transforma $X_n$ en $S_n$, la idea es unir mediante un segmento cada punto $x\in X_n$ con su correspondiente punto en $S_n$, que es $2\frac{x}{\|x\|}$. Definimos por tanto la siguiente aplicación:
                \Func{H}{X_n\times [0,1]}{X_n}{(x,t)}{(1-t)x + 2t\frac{x}{\|x\|}}
                Observemos que $\|x\|\neq 0$ para todo $x\in X_n$. Además, para comprobar que la aplicación $H$ está bien definida, debemos verificar que $H(x,t)\in X_n$ para todo $(x,t)\in X_n\times [0,1]$; es decir, que $\|H(x,t)\|\geq 2$ para todo $(x,t)\in X_n\times [0,1]$.
                \begin{align*}
                    \|H(x,t)\| &= \left\|(1-t)x + 2t\frac{x}{\|x\|}\right\| = \|x\| \left(1-t+\frac{2t}{\|x\|}\right) = \|x\|(1-t) + 2t  \\
                               &\geq 2(1-t) + 2t = 2 \qquad \forall x\in X_n, \quad \forall t\in [0,1]
                \end{align*}
                Por lo que $H(x,t)\in X_n$ para todo $(x,t)\in X_n\times [0,1]$, por lo que $H$ está bien definida. Es claro también que $H$ es continua, y observamos además que:
                \begin{itemize}
                    \item $H(x,0) = x\quad \forall x\in X_n$.
                    \item $H(x,1) = \dfrac{2x}{\|x\|}\in S_n \quad \forall x\in X_n$.
                    \item Si $y\in S_n$, entonces $\|y\|=2$, y por tanto:
                        \begin{equation*}
                            H(y,1) = 2\frac{y}{\|y\|} = y 
                        \end{equation*}
                \end{itemize}

                Por tanto, consideramos la siguiente aplicación:
                \Func{r}{X_n}{S_n}{x}{H(x,1) = 2\frac{x}{\|x\|}}
                
                Esta aplicación es continua, y además cumple que $r(y) = y$ para todo $y\in S_n$, por lo que $r$ es una retracción de $X_n$ sobre $S_n$. Para ver que $S_n$ es un retracto de deformación de $X_n$, basta con observar que $H$ es una homotopía entre $\id_{X_n}$ y $i\circ r$, donde $i:S_n\hookrightarrow X_n$ es la inclusión canónica.
            \item Estudiar si $X_1 \cong X_n$ cuando $n\geq 2$.

                Para cada $m\geq 1$ tenemos que $S_m=\{x\in \mathbb{R}^{m+1}:\|x\|=2\}$ es un retracto de deformación de $X_m$. Observamos además que la aplicación $f:\mathbb{S}^m\to S_m$ dada por:
                \begin{equation*}
                    f(x) = 2x
                \end{equation*}
                es un homeomorfismo, de donde deducimos que:
                \begin{equation*}
                    \pi_1(X_m) \cong \pi_1(S_m) \cong \pi_1(\mathbb{S}^m) \qquad \forall m\geq 1
                \end{equation*}
                Por tanto, para $n=1$ tenemos que:
                \begin{equation*}
                    \pi_1(X_1) \cong \pi_1(\mathbb{S}^1) \cong \mathbb{Z}
                \end{equation*}
                Por otro lado, para $n\geq 2$ tenemos que:
                \begin{equation*}
                    \pi_1(X_n)\cong \pi_1(\mathbb{S}^n) = \{1\}
                \end{equation*}
                Como $\{1\}\not\cong \mathbb{Z}$, deducimos que sus grupos fundamentales no son isomorfos, y por tanto $X_1\not\cong X_n$ para todo $n\geq 2$.
        \end{enumerate}
    \end{ejercicio}

    \begin{ejercicio}[2.5 puntos]
        Sea $X$ un conjunto con $|X|\geq 2$. Fijado $x_0\in X$ se considera la topología
        \begin{equation*}
            T = \{U\subseteq X : x_0\notin U\}\cup \{X\}
        \end{equation*}
        ¿Es $(X,T)$ contráctil?\\

        \noindent
        Sí es contráctil, pues si consideramos $H:X\times [0,1]\to X$ dada por:
        \begin{equation*}
            H(x,t) = \left\{\begin{array}{ll}
                x & \text{si\ } t\in \left[0,\nicefrac{1}{2}\right[ \\
                    x_0 & \text{si\ } t\in [\nicefrac{1}{2},1]
            \end{array}\right. 
        \end{equation*}
        Si tomamos $U\in T$:
        \begin{itemize}
            \item Si $U=X$ tenemos trivialmente que $H^{-1}(X) = X\times [0,1]$, que es un conjunto abierto.
            \item Si $U\neq X$ tenemos entonces que $x_0\notin U$, y:
                \begin{align*}
                    H^{-1}(U) &= \{(x,t)\in X\times [0,1] : H(x,t) \in U\} \\
                              &= \{(x,t)\in X\times [0,1]: (t\geq \nicefrac{1}{2}\land x_0 \in U) \lor (t<\nicefrac{1}{2}\land x\in U)\} \\
                              &=\{(x,t)\in X\times [0,1]: t<\nicefrac{1}{2}\land x\in U\} = U\times \left[0,\nicefrac{1}{2}\right[
                \end{align*}
                Con $U\times \left[0,\nicefrac{1}{2}\right[$ un conjunto abierto, como producto de dos conjuntos abiertos.
        \end{itemize}
        En definitiva, tenemos que $H$ es continua, así como que:
        \begin{itemize}
            \item $H(x,0) = x \quad \forall x\in X$.
            \item $H(x,1) = x_0\quad \forall x\in X$.
        \end{itemize}
        Por lo que $H$ es una homotopía entre $\id_X$ y la función constantemente igual a $x_0$, de donde $\{x_0\}$ es un retracto de deformación de $(X,T)$, por lo que este espacio topológico es contráctil.
    \end{ejercicio}

\end{document}
